% ============================================================
% DOCUMENTCLASS
% ============================================================
\documentclass[12pt,oneside]{book}

% ============================================================
% METADATOS
% ============================================================
\newcommand{\universidad}{Universidad de Buenos Aires (UBA)}
\newcommand{\institucion}{Pr\'actica Profesional Supervisada}
\newcommand{\materia}{Derecho de Familia}
\newcommand{\titulo}{Apuntes de Derecho de Familia}
\newcommand{\autor}{Isaac Edgar Camacho Ocampo}
\newcommand{\anio}{2026}

% ============================================================
% PAQUETES
% ============================================================
\usepackage[utf8]{inputenc}
\usepackage[T1]{fontenc}
\usepackage[spanish,es-nodecimaldot]{babel}
\usepackage[
    top=2.5cm,
    bottom=2.5cm,
    left=3cm,
    right=2.5cm,
    headheight=15pt
]{geometry}
\usepackage{amsmath}
\usepackage{amssymb}
\usepackage{mathtools}
\usepackage{graphicx}
\usepackage{float}
\usepackage{caption}
\usepackage{subcaption}
\usepackage{booktabs}
\usepackage{array}
\usepackage{longtable}
\usepackage{tabularx}
\usepackage{enumitem}
\usepackage{xcolor}
\usepackage[most]{tcolorbox}
\usepackage{fancyhdr}
\usepackage{ragged2e}
\usepackage[
    colorlinks=true,
    linkcolor=blue!50!black,
    citecolor=green!40!black,
    urlcolor=blue!60!black,
    pdftitle={\titulo},
    pdfauthor={\autor},
    pdfsubject={Apuntes universitarios},
    bookmarks=true
]{hyperref}

% ============================================================
% CONFIGURACIÓN
% ============================================================
\graphicspath{
    {./img/}
    {./graficos/}
    {./figuras/}
}
\setcounter{tocdepth}{2}
\setlength{\parindent}{0pt}
\setlength{\parskip}{0.5em}
\setlist[itemize]{
    topsep=0.3em,
    itemsep=0.2em
}
\setlist[enumerate]{
    topsep=0.3em,
    itemsep=0.2em
}
\clubpenalty=10000
\widowpenalty=10000

% ============================================================
% COMANDOS
% ============================================================
\newcommand{\abs}[1]{\left|#1\right|}
\newcommand{\norm}[1]{\left\lVert#1\right\rVert}
\newcommand{\vect}[1]{\mathbf{#1}}
\newcommand{\deriv}[2]{\frac{d#1}{d#2}}
\newcommand{\pderiv}[2]{\frac{\partial#1}{\partial#2}}

% ============================================================
% ENTORNOS / CAJAS
% ============================================================
\newtcolorbox{definition}[1]{
    enhanced,
    breakable,
    title={Definici\'on: #1},
    fonttitle=\bfseries,
    colback=blue!3,
    colframe=blue!50!black,
    boxrule=0.8pt,
    arc=2mm
}

\newtcolorbox{important}{
    enhanced,
    breakable,
    title={Importante},
    fonttitle=\bfseries,
    colback=yellow!8,
    colframe=orange!70!black,
    boxrule=0.8pt,
    arc=2mm
}

\newtcolorbox{example}[1]{
    enhanced,
    breakable,
    title={Ejemplo: #1},
    fonttitle=\bfseries,
    colback=green!4,
    colframe=green!40!black,
    boxrule=0.8pt,
    arc=2mm
}

\newtcolorbox{formula}[1]{
    enhanced,
    breakable,
    title={F\'ormula / Regla: #1},
    fonttitle=\bfseries,
    colback=purple!4,
    colframe=purple!50!black,
    boxrule=0.8pt,
    arc=2mm
}

\newtcolorbox{exercise}[1]{
    enhanced,
    breakable,
    title={Ejercicio: #1},
    fonttitle=\bfseries,
    colback=gray!5,
    colframe=gray!60!black,
    boxrule=0.8pt,
    arc=2mm
}

\newtcolorbox{note}{
    enhanced,
    breakable,
    title={Nota},
    fonttitle=\bfseries,
    colback=white,
    colframe=gray!50,
    boxrule=0.6pt,
    arc=2mm
}

% ============================================================
% CONFIGURACIÓN FINAL
% ============================================================
\pagestyle{fancy}
\fancyhf{}
\fancyhead[L]{\nouppercase{\leftmark}}
\fancyhead[R]{\materia}
\fancyfoot[C]{\thepage}
\renewcommand{\headrulewidth}{0.4pt}
\renewcommand{\footrulewidth}{0pt}
\fancypagestyle{plain}{
    \fancyhf{}
    \fancyfoot[C]{\thepage}
    \renewcommand{\headrulewidth}{0pt}
}

% ============================================================
% DOCUMENT
% ============================================================
\begin{document}

% ------------------------------------------------------------
% PORTADA
% ------------------------------------------------------------
\begin{titlepage}

\thispagestyle{empty}

\begin{center}

\vspace*{1cm}

{\large \textsc{\universidad}}

\vspace{0.2cm}

{\large \institucion}

\vspace{3cm}

{\Huge \textbf{\materia}}

\vspace{1cm}

{\LARGE \titulo}

\vspace{0.8cm}

\textit{Comprender los conceptos es el primer paso para convertir el estudio en conocimiento.}

\vspace{1.2cm}

\rule{0.70\textwidth}{0.5pt}

\vspace{0.5cm}

{\large Apuntes de estudio}

\vspace{0.5cm}

\rule{0.70\textwidth}{0.5pt}

\vfill

{\large \autor}

\vspace{0.3cm}

{\large \anio}

\end{center}

\vfill

\begin{flushleft}
\textit{Este es un modesto aporte a los estudiantes de ``Derecho de Familia'' no pretende sustituir el material oficial de la materia.}
\end{flushleft}

\end{titlepage}

% ------------------------------------------------------------
% ÍNDICE
% ------------------------------------------------------------
\tableofcontents

% ------------------------------------------------------------
% PRESENTACIÓN
% ------------------------------------------------------------
\chapter*{Presentaci\'on}
\addcontentsline{toc}{chapter}{Presentaci\'on}
\markboth{Presentaci\'on}{Presentaci\'on}

Estos apuntes corresponden a un pr\'actico supervisado de la materia Derecho de Familia, desarrollado en el marco de una pr\'actica profesional. La clase del 1 de noviembre de 2026 reconstruye una audiencia por violencia familiar: una medida de prohibici\'on de contacto est\'a por vencer y la defensor\'ia presta conformidad a su pr\'orroga, porque el se\~nor compareciente a\'un no acredit\'o el tratamiento terap\'eutico que el equipo interdisciplinario recomend\'o hace dos a\~nos.

A partir de ese hilo se trabajan cinco ejes:

\begin{itemize}
    \item el tratamiento terap\'eutico y su acreditaci\'on: qu\'e se exige y por qu\'e;
    \item la pr\'orroga y el dictamen no vinculante de la defensor\'ia frente a la decisi\'on de la jueza;
    \item los obst\'aculos del sistema de salud y las alternativas privadas o con copago;
    \item los domicilios, las notificaciones y la competencia: el centro de vida del ni\~no;
    \item la redacci\'on del acta y el an\'alisis posterior del equipo.
\end{itemize}

\begin{important}
\textbf{Utilidad profesional de la clase.} Seg\'un el apunte, la audiencia muestra el tipo de situaci\'on que se encuentra al comenzar a ejercer: tr\'amites trabados, operadores con posturas opuestas, plazos que vencen y clientes que no saben explicar qu\'e les pasa. Permite (1) distinguir una recomendaci\'on de una orden; (2) entender qu\'e significa acreditar; (3) anticipar problemas de domicilio y competencia; (4) cuidar el acta, porque lo que queda asentado condiciona lo que se puede pedir despu\'es; y (5) aprender a moverse en sistemas burocr\'aticos, documentar cada gesti\'on e insistir sin rendirse.
\end{important}

\begin{note}
Las intervenciones de la audiencia no estaban identificadas en el audio. Los roles que aparecen en estos apuntes (defensora, letrado, el se\~nor compareciente) se asignaron por contexto y son orientativos. Los fragmentos inaudibles del original no se reconstruyen. Se omiten los nombres propios y las direcciones concretas, que no son relevantes para el estudio.
\end{note}

\begin{example}{Una analog\'ia para relacionar los temas}
El apunte propone pensarlo como sacar el carnet de conducir para volver a manejar. ``Manejar'' es retomar el r\'egimen de comunicaci\'on con el hijo. Para eso hay que aprobar el curso (el tratamiento) y presentar el certificado (acreditarlo). Si la oficina no da turno, hay que buscar otra sede (efectores privados, Defensor\'ia de Ni\~nos). Hay que ir a la ventanilla que corresponde seg\'un el lugar de residencia (domicilio y competencia), porque si no se deriva a otra oficina y se empieza de nuevo. El acta es el comprobante sellado que guarda todo lo ocurrido.
\end{example}

% ============================================================
\chapter[Contexto y apertura]{Contexto y apertura de la audiencia}
% ============================================================

\section{Gestiones previas y tr\'amites pendientes}

La grabaci\'on comienza con una charla del equipo antes del ingreso del compareciente. El audio es fragmentario. Se mencionan las dificultades con los turnos y los canales de atenci\'on: esperas de una o dos horas y comunicaciones que se cortan, y la imposibilidad de usar el asistente virtual porque estaba asociado a otra cuenta. Tambi\'en se menciona que, respecto de un pedido, ya se hab\'ia solicitado un reiteratorio.

\begin{definition}{Reiteratorio}
Nuevo oficio que se env\'ia cuando el organismo requerido no contest\'o un oficio anterior. En la audiencia se mencionar\'a que el se\~nor ya tiene al menos tres oficios sin respuesta.
\end{definition}

\begin{note}
El apunte aclara, en una nota del editor, que Boti es el asistente virtual por WhatsApp del Gobierno de la Ciudad, usado para pedir turnos y hacer consultas.
\end{note}

\section{Contenci\'on previa al compareciente}

Antes de ingresar, se conversa con el se\~nor para tranquilizarlo y explicarle el sentido de la audiencia. Los fragmentos recuperables indican que el equipo ya anticipaba que se estaba ``al borde del impedimento de contacto'', es decir, que la medida vigente estaba cerca de su vencimiento. Se le se\~nala que cuesta hablar y recuperar la confianza, y que en la audiencia \textbf{hay que demostrar} que est\'a en condiciones de ver a su hijo.

% TODO: contenido incompleto o ambiguo en las notas originales (gran parte de este tramo es inaudible)

\begin{important}
Antes de comparecer, el se\~nor necesita confianza y claridad sobre el objetivo: la audiencia es la oportunidad de demostrar que est\'a en condiciones de retomar el v\'inculo con su hijo.
\end{important}

\section{Apertura de la audiencia}

Tras el saludo e ingreso de las partes, la primera pregunta no versa sobre el v\'inculo con el ni\~no sino sobre el tratamiento: se le pregunta cu\'ando comienza, qu\'e comenz\'o y se lo invita a contar lo hecho.

\begin{important}
El tratamiento es el punto pendiente que condiciona todo lo dem\'as; por eso es el primer tema de la audiencia.
\end{important}

\section{Cuadro Cornell: contexto y apertura}

\begin{longtable}{
    >{\raggedright\arraybackslash}p{0.28\textwidth}
    >{\raggedright\arraybackslash}p{0.66\textwidth}
}
\toprule
\textbf{Preguntas / palabras clave} & \textbf{Notas / respuestas} \\
\midrule
\endfirsthead
\toprule
\textbf{Preguntas / palabras clave} & \textbf{Notas / respuestas} \\
\midrule
\endhead
\bottomrule
\endfoot

\textbf{\textquestiondown Qu\'e es un reiteratorio?} &
Es un nuevo oficio que se env\'ia cuando el organismo requerido no contest\'o uno anterior. En el caso, el se\~nor ya ten\'ia al menos tres oficios sin respuesta. \\[0.6em]

\textbf{\textquestiondown C\'omo se prepara a quien va a comparecer?} &
Se le brinda contenci\'on, se le explica que la audiencia es la oportunidad de demostrar que est\'a en condiciones de ver a su hijo y se le recuerda que la medida vigente est\'a cerca de vencer. \\[0.6em]

\textbf{\textquestiondown Cu\'al es la primera pregunta de la audiencia y por qu\'e?} &
Se pregunta por el tratamiento (cu\'ando empieza, qu\'e empez\'o), porque es el punto pendiente que condiciona todo lo dem\'as. \\

\midrule
\multicolumn{2}{p{0.96\textwidth}}{%
\textbf{Resumen:} Antes de la audiencia se trabaj\'o la contenci\'on del compareciente y los tr\'amites pendientes (turnos, oficios y reiteratorio). La audiencia se abre con el tratamiento, porque es el asunto pendiente del que depende el resto.
} \\
\end{longtable}

% ============================================================
\chapter[Tratamiento y pr\'orroga]{Tratamiento terap\'eutico y pr\'orroga}
% ============================================================

\section{Situaci\'on del tratamiento del se\~nor}

\subsection{Gestiones realizadas}

El se\~nor relata que se present\'o el d\'ia 13 del mes en el hospital Borda con la orden para iniciar el tratamiento, pero a\'un no le asignaron turno. Indica que eligi\'o ese hospital por cercan\'ia y porque es m\'as espec\'ifico. El letrado agrega que fue varias veces y que le indicaron llamar al 147; el se\~nor llam\'o, obtuvo un turno y se acerc\'o, pero no tiene historia cl\'inica en el Borda.

\subsection{Antecedentes de atenci\'on y curso previo}

Antes, el se\~nor concurri\'o una sola vez al hospital Fern\'andez, a una cita que no fue de psicolog\'ia. No tuvo tratamiento psiqui\'atrico. Realiz\'o adem\'as un curso virtual de masculinidades, que est\'a acreditado, pero que es distinto del tratamiento terap\'eutico que se le exige. La defensora indica que buscar\'a otro sector donde gestionar un tratamiento individual.

\begin{table}[H]
\centering
\small
\begin{tabularx}{\textwidth}{>{\raggedright\arraybackslash}p{0.22\textwidth} >{\raggedright\arraybackslash}X >{\raggedright\arraybackslash}p{0.22\textwidth}}
\toprule
\textbf{Dispositivo} & \textbf{Qu\'e ocurri\'o seg\'un el se\~nor} & \textbf{Estado} \\
\midrule
Hospital Borda & Se present\'o el d\'ia 13; le indicaron llamar al 147; consigui\'o turno pero no tiene historia cl\'inica all\'i & Sin turno asignado a la fecha \\
Hospital Fern\'andez & Una \'unica consulta, no como psicolog\'ia & Sin tratamiento \\
Curso virtual de masculinidades & Lo realiz\'o & Acreditado, pero distinto del tratamiento individual \\
\bottomrule
\end{tabularx}
\end{table}

\begin{important}
El curso virtual de masculinidades no reemplaza el tratamiento terap\'eutico individual que pide el equipo interdisciplinario: son dos exigencias distintas.
\end{important}

\section{Recomendaciones del equipo interdisciplinario}

\subsection{Motivo del tratamiento}

El letrado manifiesta que no encuentra con claridad los motivos del tratamiento. La defensora responde que los motivos son las recomendaciones efectuadas en el marco de esta situaci\'on de violencia por el equipo interdisciplinario, que propuso insertar al se\~nor en un tratamiento terap\'eutico. Si esa evaluaci\'on no se cuestion\'o en su momento (cosa que pod\'ia hacerse), los profesionales pueden validar la recomendaci\'on: se trata de un cuerpo especializado al que se da intervenci\'on ante determinadas situaciones.

\subsection{Acreditar no es declarar}

La defensora explica que el tratamiento debe acreditarse y que no basta con afirmar ``hago'': hay que mostrar que se est\'a trabajando. Lo ocurrido, que gener\'o estas actuaciones, no se borra de un momento a otro. Se supone que el se\~nor desea recuperar el v\'inculo con su hijo y, en alg\'un momento, dialogar con la madre para coordinar cuestiones del ni\~no.

El letrado insiste en preguntar por qu\'e se le pide el tratamiento. La defensora reconoce que la objeci\'on tiene un punto, pues esa discusi\'on deb\'io darse en su momento.

\begin{definition}{Acreditar el tratamiento}
Seg\'un la clase, acreditar no es declarar que se hace un tratamiento, sino demostrar que se est\'a trabajando en \'el, mediante turnos, constancias y presentaciones concretas.
\end{definition}

\begin{important}
Las recomendaciones del equipo interdisciplinario emitidas en el marco de una situaci\'on de violencia no se desconocen: si no fueron cuestionadas en su momento, el tratamiento debe acreditarse, y acreditar significa demostrar un trabajo en curso, no una mera manifestaci\'on de voluntad.
\end{important}

\begin{note}
Complemento editorial del apunte (no citado en clase): se se\~nalan como marco de referencia la Ley 26.485, de Protecci\'on Integral para Prevenir, Sancionar y Erradicar la Violencia contra las Mujeres (medidas preventivas urgentes, entre ellas la prohibici\'on de acercamiento), y, en la Ciudad de Buenos Aires, la Ley 1.265 de Violencia Dom\'estica. El propio apunte indica verificar el texto vigente antes de citar.
\end{note}

\section{Conducta del se\~nor y fundamentos de la pr\'orroga}

\subsection{Lo que verifica el expediente y lo que informa el se\~nor}

El letrado sostiene que, entre los varios expedientes de violencia, no hay un hecho conductor, y que el se\~nor mantuvo di\'alogo con la madre del ni\~no \'ultimamente. Como qued\'o sin trabajo, le resulta dif\'icil pasar alimentos: aporta lo que puede, cuando puede, y se acerc\'o a entregar mercader\'ia; se le aconsej\'o que lo mejor es transferir dinero. A\'un con la orden sin renovar en distintos momentos, el se\~nor la respet\'o y se mantuvo bien desde junio. Tambi\'en se mencionan hechos nuevos y que la madre no se presenta, sin que se sepa con certeza si se present\'o o no.

% TODO: contenido incompleto o ambiguo en las notas originales (sobre la presentaci\'on de la madre)

\subsection{Por qu\'e la defensor\'ia acompa\~na la pr\'orroga}

La defensor\'ia acompa\~na la pr\'orroga de la prohibici\'on respecto del ni\~no, pero no por tener algo en contra del se\~nor: existe una recomendaci\'on de larga data. Hay un r\'egimen de comunicaci\'on iniciado que no avanzar\'a si el se\~nor no acredita que realiza una terapia. De lo contrario, \textquestiondown para qu\'e se da intervenci\'on a un equipo especializado en violencia si sus recomendaciones se omiten? La defensora no ve gravedad en ello: considera sencillo que el se\~nor se incorpore a un tratamiento para poder pedir la vinculaci\'on con su hijo y luego fijar un r\'egimen de comunicaci\'on, una vez que se hable con los terapeutas y se garantice que todo est\'a bien.

\begin{example}{El razonamiento de la defensor\'ia}
La prohibici\'on de contacto se prorroga no como sanci\'on al se\~nor, sino porque la condici\'on fijada (acreditar tratamiento) no se cumpli\'o. El r\'egimen de comunicaci\'on ya iniciado queda frenado mientras esa condici\'on siga pendiente: si no se acredita, el proceso no se mueve.
\end{example}

\section{Acceso a la salud y dictamen no vinculante}

\subsection{Los obst\'aculos del sistema de salud}

El letrado sostiene que no es lo mismo ingresar al sistema para el se\~nor que para quienes tienen otros recursos, como una prepaga: cada vez que \'el acude a la salud del Gobierno de la Ciudad, ``rebota''. Hay al menos tres oficios sin contestar, y el se\~nor no carece de voluntad.

La defensora lo entiende, pero recuerda que las recomendaciones tienen dos a\~nos, que conoce el estado del sistema de salud y que en las audiencias otras personas manifiestan el mismo inconveniente pero lograron obtener el recurso. Desde la defensor\'ia se puede oficiar y colaborar, aunque tambi\'en se perdi\'o tiempo m\'as all\'a de la complejidad de conseguir un turno. El letrado reconoce que ello es cierto.

\begin{table}[H]
\centering
\small
\begin{tabularx}{\textwidth}{>{\raggedright\arraybackslash}p{0.19\textwidth} >{\raggedright\arraybackslash}X >{\raggedright\arraybackslash}X}
\toprule
\textbf{Tema} & \textbf{Postura de la defensa (letrado)} & \textbf{Postura de la defensor\'ia} \\
\midrule
Necesidad del tratamiento & No encuentra con claridad los motivos del pedido & Surgen de las recomendaciones del equipo interdisciplinario, no cuestionadas a tiempo \\
Cumplimiento del se\~nor & Respet\'o la prohibici\'on aun sin orden renovada; mantiene di\'alogo con la madre; ayuda con mercader\'ia & Valora el esfuerzo, pero el tratamiento sigue sin acreditarse \\
Acceso a la salud & El sistema ``rebota''; hay tres oficios sin respuesta; el se\~nor no tiene prepaga & Hay personas en la misma situaci\'on que lograron turno; en dos a\~nos pudo hacerse algo m\'as \\
Salida propuesta & Comprender las dificultades reales del se\~nor & Oficiar y colaborar, pero acreditar lo pendiente \\
\bottomrule
\end{tabularx}
\end{table}

\subsection{El dictamen de la defensor\'ia no es vinculante}

La defensora aclara que, ante el pedido de pr\'orroga, el dictamen de la defensor\'ia no es vinculante: resolver\'a la jueza. No adelanta opini\'on, sino que presta conformidad con la pr\'orroga porque lo solicitado (que se acredite) no se acredit\'o. No puede decir que no se acredite y que todo est\'a bien: el tratamiento debe acreditarse y luego se dar\'a impulso al r\'egimen de comunicaci\'on, que se informa como iniciado y sin movimiento.

\begin{important}
El dictamen de la defensor\'ia no vincula a la jueza, que es quien decide. La defensor\'ia presta conformidad con la pr\'orroga porque lo que se pidi\'o (acreditar el tratamiento) no se cumpli\'o; sin esa acreditaci\'on no puede aconsejar otra cosa.
\end{important}

\section{Cuadro Cornell: tratamiento y pr\'orroga}

\begin{longtable}{
    >{\raggedright\arraybackslash}p{0.28\textwidth}
    >{\raggedright\arraybackslash}p{0.66\textwidth}
}
\toprule
\textbf{Preguntas / palabras clave} & \textbf{Notas / respuestas} \\
\midrule
\endfirsthead
\toprule
\textbf{Preguntas / palabras clave} & \textbf{Notas / respuestas} \\
\midrule
\endhead
\bottomrule
\endfoot

\textbf{\textquestiondown Qu\'e hizo el se\~nor para iniciar un tratamiento?} &
Se present\'o en el Borda el d\'ia 13 y le indicaron llamar al 147; consigui\'o turno pero no tiene historia cl\'inica all\'i y a\'un no le asignaron turno efectivo. Concurri\'o una sola vez al Fern\'andez, sin tratamiento. \\[0.6em]

\textbf{\textquestiondown El curso de masculinidades reemplaza el tratamiento?} &
No. El curso virtual est\'a acreditado, pero es distinto del tratamiento terap\'eutico individual que se exige. \\[0.6em]

\textbf{\textquestiondown Por qu\'e se le exige un tratamiento terap\'eutico?} &
Porque el equipo interdisciplinario lo recomend\'o en el marco de la situaci\'on de violencia. Si esa evaluaci\'on no se cuestion\'o en su momento, los profesionales pueden validarla: es un cuerpo especializado y recomienda lo que recomienda. \\[0.6em]

\textbf{\textquestiondown Qu\'e significa acreditar el tratamiento?} &
No basta con decir ``hago'': hay que demostrar ``estoy trabajando''. Lo ocurrido no se borra de un momento para el otro. Se acredita con turnos, constancias y presentaciones concretas. \\[0.6em]

\textbf{\textquestiondown Por qu\'e la defensor\'ia acompa\~na la pr\'orroga?} &
No por estar en contra del se\~nor, sino porque lo que se pidi\'o (acreditar) no se acredit\'o. El r\'egimen de comunicaci\'on iniciado no avanzar\'a sin terapia; omitir las recomendaciones har\'ia in\'util la intervenci\'on del equipo especializado. \\[0.6em]

\textbf{\textquestiondown C\'omo pesa la dificultad de acceso a la salud frente al dictamen?} &
La defensor\'ia reconoce las trabas (tres oficios sin respuesta, el sistema ``rebota'', el se\~nor no tiene prepaga), pero se\~nala que las recomendaciones tienen dos a\~nos y que otras personas s\'i consiguieron turno. Ofrece oficiar y colaborar, pero pide acreditar lo pendiente. \\[0.6em]

\textbf{\textquestiondown Qu\'e valor tiene el dictamen de la defensor\'ia?} &
Es no vinculante: resuelve la jueza. La defensor\'ia presta conformidad con la pr\'orroga, pero no adelanta una decisi\'on judicial. \\

\midrule
\multicolumn{2}{p{0.96\textwidth}}{%
\textbf{Resumen:} El se\~nor debe acreditar un tratamiento terap\'eutico individual recomendado hace dos a\~nos; el curso virtual no lo sustituye. La defensor\'ia presta conformidad a la pr\'orroga porque la condici\'on no se cumpli\'o, reconociendo las dificultades del sistema de salud, aunque su dictamen no vincula a la jueza. Primero se acredita y despu\'es se pide.
} \\
\end{longtable}

% ============================================================
\chapter[Domicilio y competencia]{Domicilio, notificaci\'on y competencia}
% ============================================================

\section{R\'egimen de comunicaci\'on sin movimiento}

\subsection{Por qu\'e el expediente no avanza}

La defensora explica que el r\'egimen de comunicaci\'on no tiene movimiento porque resulta muy dif\'icil notificar a la se\~nora, que ahora estar\'ia viviendo en La Matanza. El letrado observa que ahora hay un domicilio en Capital, que ser\'ia el mismo donde ella vive.

\subsection{La relaci\'on con la madre del ni\~no, seg\'un el se\~nor}

La defensora recuerda que el se\~nor dej\'o comida y vio a la madre ``hace una semana atr\'as'', como hab\'ian hablado. El se\~nor responde que se qued\'o sin trabajo, que consigui\'o una ``changuita'', que se comunic\'o con ella y que le compr\'o carne, en el lugar donde viv\'ian en pareja.

\begin{important}
El se\~nor relata que llev\'o mercader\'ia y vio a la madre. Este dato debe manejarse con cuidado cuando existe una prohibici\'on de contacto vigente o a punto de vencer: lo que el compareciente cuenta en la audiencia queda registrado. Por eso la defensor\'ia le aconseja transferir dinero en lugar de acercarse.
\end{important}

\section{Domicilios y notificaci\'on fallida}

\subsection{D\'onde vive cada uno}

El se\~nor indica el domicilio donde viv\'ian en pareja (direcci\'on a verificar). % TODO: contenido incompleto o ambiguo en las notas originales (domicilio a verificar)
Afirma que solo alquilaban all\'i y que nunca figuraron en ese domicilio; el que la se\~nora denuncia ahora es el de su DNI, que es la casa del padre, donde \'el alquila, y ella no vive all\'i. La defensora responde que ah\'i se intent\'o notificar y volvi\'o negativo. El letrado se pregunta si volvi\'o negativo porque no se encontr\'o la puerta o porque es un edificio de tres pisos con cuatro departamentos, o porque el padre trabaja todo el d\'ia.

\begin{table}[H]
\centering
\small
\begin{tabularx}{\textwidth}{>{\raggedright\arraybackslash}p{0.26\textwidth} >{\raggedright\arraybackslash}p{0.17\textwidth} >{\raggedright\arraybackslash}X}
\toprule
\textbf{Dato de domicilio} & \textbf{Qui\'en lo se\~nala} & \textbf{Observaci\'on} \\
\midrule
Domicilio denunciado en el expediente & Surge de las actuaciones & Coincide con el DNI de la se\~nora; casa de su padre, edificio de tres pisos \\
Notificaci\'on en ese domicilio & Defensor\'ia / expediente & Volvi\'o negativa; no queda claro si por no hallar la puerta o por ausencia del padre \\
Domicilio donde viv\'ia la pareja & El se\~nor & Era alquilado; el se\~nor afirma que la se\~nora vive all\'i actualmente \\
Residencia en provincia & Defensor\'ia & Se menciona La Matanza (Lomas del Mirador) como posible residencia de la madre y del ni\~no \\
\bottomrule
\end{tabularx}
\end{table}

\subsection{Primer ofrecimiento de un listado de efectores privados}

La defensora, para no perder tiempo, ofrece algunos lugares donde consultar por el tratamiento, dado que el 147 y el Borda est\'an colapsados. Son lugares privados, algunos con copago, pero mucho m\'as accesibles que otras opciones privadas, y trabajan con varones.

\subsection{Primer pedido de constancia en el acta}

La defensora solicita que se deje asentado en el acta que el se\~nor dice que el domicilio denunciado es el del padre de la se\~nora, y que la se\~nora vive actualmente en el domicilio anterior, donde \'el le dej\'o la carne hace tres semanas; pide tambi\'en que se consigne la direcci\'on exacta de ese domicilio.

\section{Centro de vida del ni\~no y competencia}

\subsection{D\'onde est\'a el ni\~no}

La defensora pregunta si el ni\~no va al jard\'in en provincia; el se\~nor dice que no lo sabe. Se se\~nala que se desconoce si la se\~nora trabaja en Capital y que no se comprende por qu\'e las actuaciones se tramitaron en Capital si viv\'ian en otro lugar. Se entiende que la denuncia se realiz\'o en la OVD de Capital y que en 2024 la se\~nora viv\'ia en Capital, en la casa de su padre.

\subsection{El recorrido de domicilios del ni\~no}

Seg\'un el se\~nor, el ni\~no nunca vivi\'o en Capital: al nacer vivieron un tiempo en la casa de los padres de \'el, un tiempo en la casa del padre de ella, y despu\'es se fueron a Lomas del Mirador, donde la madre qued\'o. El se\~nor sigue viviendo en Lugano. El ni\~no concurri\'o a un CPI en Mataderos y hoy ve a un psic\'ologo en Capital, lo que es m\'as reciente. La medida se dispuso en Capital y la notificaci\'on se intent\'o all\'i; se advirti\'o la situaci\'on al insistir el se\~nor en llevar comida al ni\~no.

\begin{example}{Un mismo ni\~no, varios domicilios}
Seg\'un el relato: al nacer, el ni\~no vivi\'o un tiempo en casa de los padres del se\~nor y luego en casa del padre de la se\~nora; despu\'es la madre pas\'o a vivir en Lomas del Mirador (provincia). El ni\~no asisti\'o a un CPI en Mataderos y hoy ve a un psic\'ologo en Capital. Estas idas y venidas hacen que la denuncia se haya tramitado en Capital aunque el ni\~no pueda tener su centro de vida en provincia.
\end{example}

\subsection{La competencia}

La defensora sostiene que no es un dato menor: el centro de vida del ni\~no es donde debe estar la cercan\'ia. Si el r\'egimen de comunicaci\'on ya est\'a iniciado pero el domicilio del ni\~no es en provincia, no tramitar\'a en Capital. Aclara que la evaluaci\'on se hizo y nadie cuestion\'o esto en su momento, porque el ni\~no asist\'ia a un CPI en Capital.

\begin{definition}{Centro de vida}
Complemento editorial del apunte (no citado en clase): la Ley 26.061, art. 3, inc. f, define el centro de vida como el lugar donde ni\~nas, ni\~nos y adolescentes hubiesen transcurrido, en condiciones leg\'itimas, la mayor parte de su existencia. Seg\'un el mismo complemento, el art. 716 del C\'odigo Civil y Comercial contiene una regla de competencia en los procesos de responsabilidad parental vinculada con el centro de vida; el apunte indica verificar el texto vigente antes de citar.
\end{definition}

\begin{important}
Si el ni\~no vive en provincia, el juzgado de Capital puede declararse incompetente y remitir las actuaciones. La defensora pide que quede asentado en el acta y quiere saber d\'onde est\'a el ni\~no antes de que empiece primer grado.
\end{important}

\section{Cuadro Cornell: domicilio y competencia}

\begin{longtable}{
    >{\raggedright\arraybackslash}p{0.28\textwidth}
    >{\raggedright\arraybackslash}p{0.66\textwidth}
}
\toprule
\textbf{Preguntas / palabras clave} & \textbf{Notas / respuestas} \\
\midrule
\endfirsthead
\toprule
\textbf{Preguntas / palabras clave} & \textbf{Notas / respuestas} \\
\midrule
\endhead
\bottomrule
\endfoot

\textbf{\textquestiondown Por qu\'e el r\'egimen de comunicaci\'on no avanza?} &
Porque resulta muy dif\'icil notificar a la se\~nora, que ahora vivir\'ia en La Matanza. El r\'egimen est\'a iniciado pero ``sin movimiento''. \\[0.6em]

\textbf{\textquestiondown D\'onde vive realmente cada parte?} &
El domicilio denunciado es el del DNI (casa del padre de la se\~nora) y all\'i la notificaci\'on volvi\'o negativa. El se\~nor sostiene que ella vive en el antiguo domicilio de la pareja. Tambi\'en se menciona La Matanza como posible residencia de madre e hijo. \\[0.6em]

\textbf{\textquestiondown Por qu\'e es delicado el contacto del se\~nor con la madre?} &
Con una prohibici\'on vigente o por vencer, lo que el compareciente cuenta en la audiencia queda registrado; por eso se le aconseja transferir dinero en lugar de acercarse. \\[0.6em]

\textbf{\textquestiondown Qu\'e es el centro de vida y por qu\'e importa?} &
La defensora lo se\~nala como dato central: all\'i debe estar la cercan\'ia. El apunte agrega, como complemento editorial no citado en clase, la Ley 26.061, art. 3, inc. f. \\[0.6em]

\textbf{\textquestiondown Qu\'e juzgado es competente seg\'un d\'onde vive el ni\~no?} &
Si el ni\~no vive en provincia (Lomas del Mirador), el juzgado de Capital puede declararse incompetente y remitir el caso. La defensora quiere saber d\'onde est\'a el ni\~no antes de que comience primer grado. \\

\midrule
\multicolumn{2}{p{0.96\textwidth}}{%
\textbf{Resumen:} Los domicilios cruzados y la notificaci\'on fallida explican que el r\'egimen de comunicaci\'on est\'e sin movimiento. Adem\'as, si el centro de vida del ni\~no est\'a en provincia, el juzgado de Capital podr\'ia declararse incompetente, por lo que la defensora pide dejar constancia y averiguar d\'onde est\'a el ni\~no.
} \\
\end{longtable}

% ============================================================
\chapter[Cierre: acta y an\'alisis]{Cierre: acta y an\'alisis posterior}
% ============================================================

\section{Alternativas de tratamiento y consejos}

\subsection{El consejo central}

La defensora aconseja que, con el listado de efectores privados, el se\~nor acredite el inicio del tratamiento. El se\~nor advierte que hay mucha violencia y que trabajan con ni\~nos; la defensora le indica que se fije sobre todo en los que trabajan con varones, que ella fue anotando.

\subsection{La Defensor\'ia de Ni\~nos de la Comuna 9}

El se\~nor cuenta que recibi\'o un mensaje de la Defensor\'ia de Ni\~nos de la Comuna 9, que lo convoc\'o (se trata tambi\'en de un equipo interdisciplinario). Ten\'ia cita para el jueves 22, pero oy\'o el mensaje al d\'ia siguiente. Dej\'o su n\'umero por mail y por tel\'efono, sin respuesta. La defensora le aconseja insistir o concurrir, porque a veces est\'an desbordados de casos, y recuerda que se trata de un equipo con el que tambi\'en trabajan la jueza y la defensor\'ia, y que puede ser un puente para derivarlo a un lugar donde hacer terapia. La preocupaci\'on del se\~nor, se\~nala, es ir armando de a poco el v\'inculo con su hijo.

\subsection{C\'omo se ordena el pr\'oximo paso}

La defensora indica que acrediten cuanto antes que se movieron, que tienen turno, y que despu\'es podr\'an pedir lo que quieran. Luego deber\'a evaluarse la competencia, que quiere que quede asentada en el acta, pues no es lo mismo la urgencia de una situaci\'on de violencia que el ni\~no empiece primer grado en provincia o viaje todos los d\'ias a Capital. Es importante acreditar lo pendiente para poder avanzar en cualquier otra petici\'on.

\begin{important}
\textbf{Primero acreditar, despu\'es pedir.} Mientras el tratamiento no est\'e acreditado, cualquier otra petici\'on del se\~nor (como avanzar con el r\'egimen de comunicaci\'on) queda condicionada. La Defensor\'ia de Ni\~nos de la Comuna 9 puede funcionar adem\'as como puente de derivaci\'on a un espacio terap\'eutico.
\end{important}

\section{Redacci\'on del acta}

\subsection{Qu\'e se asienta y c\'omo}

Durante la redacci\'on se confirma el domicilio de la se\~nora, donde tendr\'ia un local. La defensora indica que se consigne la entrega del listado y el compromiso del se\~nor de acreditar: si ya tiene turno, lo presenta y luego se ver\'a. Tambi\'en que el se\~nor se ofici\'o tres veces y que present\'o el d\'ia 13 en el Borda el papel para iniciar tratamiento, sin turno asignado hasta el momento, por lo que debe buscarse otro efector. Se menciona una constancia que acredita que efectu\'o la presentaci\'on en el Borda sin haber obtenido respuesta a la fecha.

\subsection{Constancia sobre el domicilio y la convocatoria}

Se asienta que el domicilio aportado por la se\~nora, a la fecha, es la casa de su familia paterna. % TODO: contenido incompleto o ambiguo en las notas originales (dato a verificar sobre el domicilio)
Se deja constancia de que el se\~nor fue convocado ante la Defensor\'ia de la Comuna 9, aunque probablemente la fecha sea tard\'ia, y que se reiterar\'a ante esa defensor\'ia para pedir una nueva fecha de convocatoria.

\subsection{Cierre y firmas}

No se piden m\'as constancias; luego el se\~nor ver\'a qu\'e se resuelve ante los pedidos efectuados. Se se\~nala que lo pendiente est\'a acreditado en el expediente y que la intimaci\'on ten\'ia que ver con el otro tratamiento. Se asientan los datos del compareciente (inaudibles en el audio) y se procede a las firmas, con la aclaraci\'on debajo de cada una.

\begin{important}
Quedan asentadas en el acta: (1) la entrega del listado de efectores y el compromiso del se\~nor de acreditar; (2) los tres oficios sin respuesta; (3) la constancia de la presentaci\'on en el Borda sin respuesta; (4) la versi\'on del domicilio (el denunciado es el del padre de la se\~nora; el actual ser\'ia el anterior domicilio de la pareja); y (5) la convocatoria ante la Defensor\'ia de Ni\~nos y el pedido de nueva fecha.
\end{important}

\section{An\'alisis posterior del equipo}

Una vez que el se\~nor se retira, la docente y los practicantes comentan lo sucedido.

\subsection{Primeras impresiones}

Un practicante observa que la evaluaci\'on es de hace mucho tiempo y que, si el se\~nor no se comprometi\'o, son dos a\~nos sin avance.

\subsection{Recomendaci\'on frente a orden}

Se comenta que al se\~nor se le dec\'ia que lo suyo era una recomendaci\'on y no una orden. La docente afirma que las recomendaciones vienen con una orden por detr\'as.

\begin{table}[H]
\centering
\small
\begin{tabularx}{\textwidth}{>{\raggedright\arraybackslash}p{0.2\textwidth} >{\raggedright\arraybackslash}X}
\toprule
\textbf{Concepto} & \textbf{Lectura transmitida en la clase} \\
\midrule
Recomendaci\'on & Aparenta ser opcional, pero proviene de un cuerpo especializado con intervenci\'on judicial; en la pr\'actica condiciona el avance del proceso \\
Orden & Obligaci\'on expresa; si no se cumple, hay consecuencias procesales directas \\
Relaci\'on entre ambas & La recomendaci\'on ``viene con una orden por detr\'as'': si no se acredita, el r\'egimen de comunicaci\'on no avanza \\
\bottomrule
\end{tabularx}
\end{table}

\subsection{Plan de acci\'on del equipo}

La docente quiere ver qu\'e se sube al acta, y se\~nala que no se acreditar\'a hasta que se acredite lo pendiente. Se informa que ya se est\'a notificando el domicilio de Lomas del Mirador. La docente advierte que, si el ni\~no efectivamente vive all\'i, el juzgado se declarar\'a incompetente y remitir\'a las actuaciones, y ``ah\'i ya no hacemos nada''. Agrega que, si el se\~nor no hace terapia, no podr\'a avanzar con lo que quiere.

\subsection{La lectura sobre el compareciente: miedo, silencio y timidez}

La docente y los practicantes comentan que el se\~nor parece muy t\'imido, con dificultad para hablar y expresarse; en un momento se quebr\'o y le salieron l\'agrimas. Se plantea acompa\~narlo a gestionar el turno en el Borda, con la cautela de no sobrecargar a\'un m\'as un sistema colapsado.

\begin{note}
El apunte se\~nala que la \'ultima l\'inea del archivo original (``Sincronizar con el tiempo del video'') es un elemento de la interfaz de la plataforma de video y no forma parte de la clase; fue omitida.
\end{note}

\section{Cuadro Cornell: cierre, acta y an\'alisis}

\begin{longtable}{
    >{\raggedright\arraybackslash}p{0.28\textwidth}
    >{\raggedright\arraybackslash}p{0.66\textwidth}
}
\toprule
\textbf{Preguntas / palabras clave} & \textbf{Notas / respuestas} \\
\midrule
\endfirsthead
\toprule
\textbf{Preguntas / palabras clave} & \textbf{Notas / respuestas} \\
\midrule
\endhead
\bottomrule
\endfoot

\textbf{\textquestiondown Qu\'e alternativas concretas tiene el se\~nor?} &
Un listado de efectores privados con copago y que trabajan con varones; reiterar ante la Defensor\'ia de Ni\~nos de la Comuna 9 (puente de derivaci\'on); presentar el turno apenas exista. \\[0.6em]

\textbf{\textquestiondown Qu\'e significa ``primero acreditar, despu\'es pedir''?} &
Mientras el tratamiento no est\'e acreditado, cualquier otra petici\'on, como avanzar con el r\'egimen de comunicaci\'on, queda condicionada. \\[0.6em]

\textbf{\textquestiondown Qu\'e debe quedar asentado en el acta?} &
Entrega del listado y compromiso de acreditar; tres oficios; constancia del Borda sin respuesta; versi\'on de los domicilios; convocatoria ante la Defensor\'ia de Ni\~nos y pedido de nueva fecha. \\[0.6em]

\textbf{\textquestiondown Qu\'e diferencia hay entre recomendaci\'on y orden?} &
La recomendaci\'on parece optativa, pero ``viene con una orden por detr\'as'': si no se acredita, el r\'egimen de comunicaci\'on no avanza. Su valor pr\'actico es el de una condici\'on. \\[0.6em]

\textbf{\textquestiondown Qu\'e lecturas profesionales deja la audiencia?} &
Se percibe al compareciente muy t\'imido, con dificultad para expresarse; se sugiere acompa\~narlo en las gestiones, sin sobrecargar un sistema ya colapsado. Si el ni\~no vive en provincia, el juzgado puede declararse incompetente. \\

\midrule
\multicolumn{2}{p{0.96\textwidth}}{%
\textbf{Resumen:} El acta fija las constancias, el listado de efectores y el compromiso del se\~nor. El an\'alisis posterior muestra que una ``recomendaci\'on'' profesional tiene en la pr\'actica fuerza de orden, y que acompa\~nar a una persona t\'imida o con miedo es parte del trabajo del abogado.
} \\
\end{longtable}

% ============================================================
\chapter[Repaso general]{Repaso general (Cornell)}
% ============================================================

\begin{longtable}{
    >{\raggedright\arraybackslash}p{0.28\textwidth}
    >{\raggedright\arraybackslash}p{0.66\textwidth}
}
\toprule
\textbf{Preguntas / palabras clave} & \textbf{Notas / respuestas} \\
\midrule
\endfirsthead
\toprule
\textbf{Preguntas / palabras clave} & \textbf{Notas / respuestas} \\
\midrule
\endhead
\bottomrule
\endfoot

\textbf{\textquestiondown Por qu\'e se le exige tratamiento terap\'eutico al se\~nor?} &
Porque el equipo interdisciplinario lo recomend\'o en el marco de la situaci\'on de violencia. El tratamiento es condici\'on para avanzar con el v\'inculo. \\[0.6em]

\textbf{\textquestiondown Qu\'e significa acreditar?} &
Demostrar que se est\'a trabajando, con turnos, constancias y presentaciones concretas; no alcanza con decir ``hago''. \\[0.6em]

\textbf{Recomendaci\'on vs. orden} &
La recomendaci\'on ``viene con una orden por detr\'as'': si no se acredita, el r\'egimen de comunicaci\'on no avanza. \\[0.6em]

\textbf{\textquestiondown Qu\'e valor tiene el dictamen de la defensor\'ia?} &
No es vinculante: resuelve la jueza. La defensor\'ia presta conformidad con la pr\'orroga porque no se acredit\'o lo pedido. \\[0.6em]

\textbf{Dificultades del sistema de salud} &
Tres oficios sin respuesta; el 147 y el Borda no dieron turno. La defensor\'ia reconoce las trabas, pero se\~nala que las recomendaciones tienen dos a\~nos y que otras personas consiguieron turno. \\[0.6em]

\textbf{Domicilio y notificaci\'on} &
El domicilio denunciado es el del DNI (casa del padre de la se\~nora) y all\'i la notificaci\'on volvi\'o negativa; esto explica que el r\'egimen est\'e ``sin movimiento''. \\[0.6em]

\textbf{Centro de vida y competencia} &
Si el ni\~no vive en provincia, el juzgado de Capital puede declararse incompetente y remitir el caso. \\

\midrule
\multicolumn{2}{p{0.96\textwidth}}{%
\textbf{Resumen:} La audiencia gira en torno a una condici\'on: el se\~nor no podr\'a retomar el r\'egimen de comunicaci\'on con su hijo mientras no acredite un tratamiento terap\'eutico recomendado hace dos a\~nos por el equipo interdisciplinario. Por eso la defensor\'ia presta conformidad a la pr\'orroga de la prohibici\'on de contacto, aun reconociendo que el sistema de salud presenta trabas reales. El dictamen de la defensor\'ia no vincula a la jueza, pero deja claro que primero se acredita y despu\'es se pide. En paralelo aparece un problema de fondo: los domicilios cruzados, la notificaci\'on fallida y la posibilidad de que el ni\~no tenga su centro de vida en provincia, lo que podr\'ia llevar a una declaraci\'on de incompetencia. El acta fija las constancias, el listado de efectores y el compromiso del se\~nor, y el an\'alisis posterior muestra que una ``recomendaci\'on'' profesional tiene en la pr\'actica fuerza de orden.
} \\
\end{longtable}

\end{document}
